% GENERATED FILE. Edit canonical lessons, then run npm run generate:books.
% canonical-source-hash: fnv1a64:2e69b4c84a8120c7
% canonical-lessons: MR-C224-mhatare-hone, MR-C224-ghatasphot-ghene, MR-C224-ghar-badalne, MR-C224-aram-karne, MR-C224-kantalne

\chapter{To grow old, To get divorced, To move house, To rest, To get bored}
\label{ch:mr-to-grow-old-to-get-divorced-to-move-house-to-rest-to-get-bored}
\hlchaptermodality{224}

\begin{chapteropening}
\textbf{By the end of this chapter:} \emph{I can say to grow old, to get divorced, to move house, to rest and to get bored.}

Complete the last lesson of chapter 224: I can say to grow old, to get divorced, to move house, to rest and to get bored.
\end{chapteropening}

\section[\texorpdfstring{\mr{म्हातारे} \mr{होणे} (mhātāre hoṇe)}{म्हातारे होणे (mhātāre hoṇe)}]{\mr{म्हातारे} \mr{होणे} (mhātāre hoṇe) — to grow old}

\label{lesson:MR-C224-mhatare-hone}



\begin{quote}
Before the new one: say the Marathi for to translate, then the Marathi for to make a decision.
\end{quote}

\subsection*{You'll want to know: \mr{म्हातारे} \mr{होणे}}
\textbf{\mr{म्हातारे} \mr{होणे}} — \emph{mhātāre hoṇe} — \textquotedblleft{}to grow old\textquotedblright{}.

\begin{grammarlens}[title={What you've built}]
One more word for this chapter.
\end{grammarlens}

\subsection*{Your turn}
\begin{itemize}
\raggedright
  \item \emph{Say it:} \emph{mhātāre hoṇe}
  \item \emph{Say it:} \emph{mhātāre hoṇe}, once more
  \item \emph{Say it:} say \emph{mhātāre hoṇe}
  \item \emph{Recall:} say the Marathi for to translate, then the Marathi for to make a decision, then say \emph{mhātāre hoṇe} again
\end{itemize}

\begin{tcolorbox}[breakable,colback=teal!4,colframe=teal!35!black,title={Before you move on}]
What does \mr{म्हातारे} \mr{होणे} mean? (\textquotedblleft{}To grow old\textquotedblright{}.) Say it once more.
\end{tcolorbox}
\section[\texorpdfstring{\mr{घटस्फोट} \mr{घेणे} (ghaṭasphoṭ gheṇe)}{घटस्फोट घेणे (ghaṭasphoṭ gheṇe)}]{\mr{घटस्फोट} \mr{घेणे} (ghaṭasphoṭ gheṇe) — to get divorced}

\label{lesson:MR-C224-ghatasphot-ghene}



\begin{quote}
Before the new one: say the Marathi for to make a decision, then the Marathi for to grow old.
\end{quote}

\subsection*{You'll want to know: \mr{घटस्फोट} \mr{घेणे}}
\textbf{\mr{घटस्फोट} \mr{घेणे}} — \emph{ghaṭasphoṭ gheṇe} — \textquotedblleft{}to get divorced\textquotedblright{}.

\begin{grammarlens}[title={What you've built}]
One more word for this chapter.
\end{grammarlens}

\subsection*{Your turn}
\begin{itemize}
\raggedright
  \item \emph{Say it:} \emph{ghaṭasphoṭ gheṇe}
  \item \emph{Say it:} \emph{ghaṭasphoṭ gheṇe}, once more
  \item \emph{Say it:} say \emph{ghaṭasphoṭ gheṇe}
  \item \emph{Recall:} say the Marathi for to make a decision, then the Marathi for to grow old, then say \emph{ghaṭasphoṭ gheṇe} again
\end{itemize}

\begin{tcolorbox}[breakable,colback=teal!4,colframe=teal!35!black,title={Before you move on}]
What does \mr{घटस्फोट} \mr{घेणे} mean? (\textquotedblleft{}To get divorced\textquotedblright{}.) Say it once more.
\end{tcolorbox}
\section[\texorpdfstring{\mr{घर} \mr{बदलणे} (ghar badalṇe)}{घर बदलणे (ghar badalṇe)}]{\mr{घर} \mr{बदलणे} (ghar badalṇe) — to move house}

\label{lesson:MR-C224-ghar-badalne}



\begin{quote}
Before the new one: say the Marathi for to grow old, then the Marathi for to get divorced.
\end{quote}

\subsection*{You'll want to know: \mr{घर} \mr{बदलणे}}
\textbf{\mr{घर} \mr{बदलणे}} — \emph{ghar badalṇe} — \textquotedblleft{}to move house\textquotedblright{}.

\begin{grammarlens}[title={What you've built}]
One more word for this chapter.
\end{grammarlens}

\subsection*{Your turn}
\begin{itemize}
\raggedright
  \item \emph{Say it:} \emph{ghar badalṇe}
  \item \emph{Say it:} \emph{ghar badalṇe}, once more
  \item \emph{Say it:} say \emph{ghar badalṇe}
  \item \emph{Recall:} say the Marathi for to grow old, then the Marathi for to get divorced, then say \emph{ghar badalṇe} again
\end{itemize}

\begin{tcolorbox}[breakable,colback=teal!4,colframe=teal!35!black,title={Before you move on}]
What does \mr{घर} \mr{बदलणे} mean? (\textquotedblleft{}To move house\textquotedblright{}.) Say it once more.
\end{tcolorbox}
\section[\texorpdfstring{\mr{आराम} \mr{करणे} (ārām karṇe)}{आराम करणे (ārām karṇe)}]{\mr{आराम} \mr{करणे} (ārām karṇe) — to rest}

\label{lesson:MR-C224-aram-karne}



\begin{quote}
Before the new one: say the Marathi for to get divorced, then the Marathi for to move house.
\end{quote}

\subsection*{You'll want to know: \mr{आराम} \mr{करणे}}
\textbf{\mr{आराम} \mr{करणे}} — \emph{ārām karṇe} — \textquotedblleft{}to rest\textquotedblright{}.

\begin{grammarlens}[title={What you've built}]
One more word for this chapter.
\end{grammarlens}

\subsection*{Your turn}
\begin{itemize}
\raggedright
  \item \emph{Say it:} \emph{ārām karṇe}
  \item \emph{Say it:} \emph{ārām karṇe}, once more
  \item \emph{Say it:} say \emph{ārām karṇe}
  \item \emph{Recall:} say the Marathi for to get divorced, then the Marathi for to move house, then say \emph{ārām karṇe} again
\end{itemize}

\begin{tcolorbox}[breakable,colback=teal!4,colframe=teal!35!black,title={Before you move on}]
What does \mr{आराम} \mr{करणे} mean? (\textquotedblleft{}To rest\textquotedblright{}.) Say it once more.
\end{tcolorbox}
\section[\texorpdfstring{\mr{कंटाळणे} (kaṇṭāḷṇe)}{कंटाळणे (kaṇṭāḷṇe)}]{\mr{कंटाळणे} (kaṇṭāḷṇe) — to get bored}

\label{lesson:MR-C224-kantalne}



\begin{quote}
Before the new one: say the Marathi for to move house, then the Marathi for to rest.
\end{quote}

\subsection*{You'll want to know: \mr{कंटाळणे}}
\textbf{\mr{कंटाळणे}} — \emph{kaṇṭāḷṇe} — \textquotedblleft{}to get bored\textquotedblright{}.

\begin{grammarlens}[title={What you've built}]
One more word for this chapter.
\end{grammarlens}

\subsection*{Your turn}
\begin{itemize}
\raggedright
  \item \emph{Say it:} \emph{kaṇṭāḷṇe}
  \item \emph{Say it:} \emph{kaṇṭāḷṇe}, once more
  \item \emph{Say it:} say \emph{kaṇṭāḷṇe}
  \item \emph{Recall:} say the Marathi for to move house, then the Marathi for to rest, then say \emph{kaṇṭāḷṇe} again
\end{itemize}

\begin{tcolorbox}[breakable,colback=teal!4,colframe=teal!35!black,title={Before you move on}]
What does \mr{कंटाळणे} mean? (\textquotedblleft{}To get bored\textquotedblright{}.) Say it once more.
\end{tcolorbox}
